\section{Finite parts, regularization and residual readouts}
\label{lectura:manuscrito-04}
The functional constructed from integer modes admits operations preserving different aspects of its structure. Evaluating it at a point of convergence, extracting the regular term accompanying a pole and differentiating its meromorphic continuation are three precise operations. This distinction organizes the constants of this chapter and their relationship with those in the moments chapter.

The arithmetic antecedent is representation of the normal forms obtained with the two APP sheets. Their evaluation identifies each positive integer once; the enumeration operator acts on the corresponding basis by

\[
Qe_n=ne_n,\qquad
\operatorname{Dom}Q=\left\{u:\sum_{n\geq1}n^2|u_n|^2<\infty\right\}.
\]
The subsequent spectral readout is

\[
Z(s)=\operatorname{Tr}(Q^{-s})=\sum_{n\geq1}n^{-s},\qquad \Re s>1.
\]
The chapter on the completed form proves its agreement with the theta–Mellin functional and its continuation. Here we bring together an evaluation procedure with controlled remainders. History memory remains in the source state; the trace is a readout of expressed modes, not a reconstruction of that entire history from a scalar constant.

\subsection{A continuation formula with explicit remainder}
Define the Bernoulli numbers through the formal identity

\[
\frac{t}{e^t-1}=\sum_{j\geq0}B_j\frac{t^j}{j!}.
\]
Their coefficients are computed without a numerical table:

\[
B_0=1,\qquad
B_m=-\frac1{m+1}\sum_{j=0}^{m-1}\binom{m+1}{j}B_j.
\]
In particular, \(B_1=-1/2\). Write \((s)_r=s(s+1)\cdots(s+r-1)\), with \((s)_0=1\), and denote the corresponding polynomial by \(B_m(x)=\sum_{j=0}^m\binom mjB_jx^{m-j}\). Its periodization is \(\widetilde B_{2p}(x)=B_{2p}(\{x\})\). For \(N,p\geq1\), Euler–Maclaurin summation gives

\[
\begin{aligned}
Z(s)={}&\sum_{n=1}^N n^{-s}+\frac{N^{1-s}}{s-1}-\frac1{2N^s}\\
&+\sum_{k=1}^p\frac{B_{2k}}{(2k)!}(s)_{2k-1}N^{1-s-2k}
+R_{N,p}(s),
\end{aligned}
\]
with

\[
R_{N,p}(s)=-\frac{(s)_{2p}}{(2p)!}
\int_N^\infty\widetilde B_{2p}(x)x^{-s-2p}\,dx.
\]
\HMTresultado{Proposition}{Domain and bound of the remainder}{rz-num-04-1}
The remainder integral is holomorphic on \(\Re s>1-2p\). If \(\sigma=\Re s>1-2p\), then

\[
|R_{N,p}(s)|\leq
\frac{|B_{2p}|\,|(s)_{2p}|}{(2p)!}
\frac{N^{1-\sigma-2p}}{\sigma+2p-1}.
\]
The preceding formula continues \(Z\) to that half-plane, with its simple pole at \(s=1\), and is compatible with the theta–Mellin continuation.

\textbf{Proof.} The summation formula is obtained by successive integration by parts on each unit interval, using \(B_j'(x)=jB_{j-1}(x)\) and the endpoint equalities for polynomials of degree other than one. Applying it to \(x^{-s}\) on \([N,M]\), the derivatives supply the factors \((s)_r\). On \(\Re s>1\), passage to \(M\to\infty\) removes the upper endpoint terms and produces exactly the identity and remainder sign written above.

To control the remainder, use

\[
|B_{2p}(t)|\leq |B_{2p}|,\qquad 0\leq t\leq1.
\]
One proof of this inequality is the absolutely convergent Fourier series

\[
B_{2p}(t)=(-1)^{p+1}\frac{2(2p)!}{(2\pi)^{2p}}
\sum_{m\geq1}\frac{\cos(2\pi mt)}{m^{2p}}.
\]
The maximum of the absolute-value bound is attained at \(t=0\). This identity follows from the Fourier coefficients by integration by parts of the polynomials themselves; the factor \(\pi\) merely expresses the Fourier convention of this proof and does not enter numerically into the rational evaluator used later. Comparison with \(x^{-\sigma-2p}\) proves the remainder bound. On compact subsets of the half-plane, it also dominates derivatives in \(s\), which add powers of \(\log x\); the integral is therefore holomorphic there. The identity on the convergence half-plane and uniqueness of continuation prove the stated compatibility. \hfill\(\square\)

Order \(p\) and horizon \(N\) are computation and error-control parameters. They do not select a different constant: expressions with different \(N,p\) represent the same functional on their common domains. Nor must an asymptotic series converge as \(p\) increases indefinitely with \(N\) fixed. Each finite pair comes with its own proved bound.

\subsection{Euler–Mascheroni as the finite part of the same functional}
The finite part at the pole is defined by

\[
\gamma_0=\lim_{s\to1}\left(Z(s)-\frac1{s-1}\right).
\]
Setting \(s=1+z\), the singular term satisfies

\[
\frac{N^{-z}}z=\frac1z-\log N+O(z).
\]
The preceding proposition allows the finite part to be taken term by term.

\HMTresultado{Theorem}{Explicit interval for the finite part}{rz-num-04-2}
Let

\[
G_{N,p}=H_N-\log N-\frac1{2N}
+\sum_{k=1}^p\frac{B_{2k}}{2kN^{2k}},
\qquad H_N=\sum_{n=1}^N\frac1n.
\]
Then

\[
|\gamma_0-G_{N,p}|\leq\frac{|B_{2p}|}{2pN^{2p}}.
\]
Moreover,

\[
\gamma_0=\lim_{N\to\infty}(H_N-\log N).
\]
\textbf{Proof.} At \(s=1\), one has \((1)_{2k-1}=(2k-1)!\); hence each Bernoulli term becomes \(B_{2k}/(2kN^{2k})\). The bound in Proposition~\ref{rz-num-04-1}, evaluated at that point, is \(|B_{2p}|/(2pN^{2p})\). With \(p=1\), all terms separating \(G_{N,1}\) from \(H_N-\log N\), and the remainder as well, tend to zero. This proves the limit. This is the subsequent identification with the Euler–Mascheroni constant. \hfill\(\square\)

The antecedent of \(\gamma_0\) used in the Meissel–Mertens formula is thus incorporated. The number is not received as an independent numerical value to correct the sum over irreducibles: it is a specified operation on the same integer partition function.

The logarithm used in the calculation also admits a rational interval. For \(x\geq1\), let \(z=(x-1)/(x+1)\). The integrated series of a geometric progression gives

\[
\log x=2\sum_{j=0}^{M-1}\frac{z^{2j+1}}{2j+1}+r_M(x),
\quad
0\leq r_M(x)\leq
\frac{2z^{2M+1}}{(2M+1)(1-z^2)}.
\]
The inequality follows by replacing the tail denominators with \(2M+1\) and summing the geometric progression. First reduce \(x=2^a y\), with \(1\leq y<2\), and evaluate \(a\log2+\log y\). The series argument is then at most \(1/3\). For \(0<x<1\), use \(\log x=-\log(1/x)\), reversing the interval endpoints.

\subsection{Apéry and the regularized negative evaluation}
For an integer \(r\geq2\), specialization of Proposition~\ref{rz-num-04-1} produces a rational centre

\[
Z_{N,p}^{[r]}=
\sum_{n=1}^N\frac1{n^r}+\frac1{(r-1)N^{r-1}}-\frac1{2N^r}
+\sum_{k=1}^p\frac{B_{2k}(r)_{2k-1}}{(2k)!N^{r+2k-1}}
\]
and the entirely rational bound

\[
\left|Z(r)-Z_{N,p}^{[r]}\right|
\leq\frac{|B_{2p}|(r)_{2p-1}}{(2p)!N^{r+2p-1}}.
\]
The simplification uses \((r)_{2p}/(r+2p-1)=(r)_{2p-1}\). At \(r=3\), one obtains Apéry's constant. The procedure determines its value with bounds; irrationality of that constant is a mathematical statement distinct from this evaluation.

On the negative side, the same continuation admits an exact evaluation, without a decimal approximation or fitting parameter.

\HMTresultado{Theorem}{Regularized evaluation at minus one}{rz-num-04-3}
The continuation of the integer-mode functional satisfies

\[
Z(-1)=-\frac1{12}.
\]
\textbf{Proof.} In the formula of Proposition~\ref{rz-num-04-1}, take \(p=2\), so that \(s=-1\) lies in the interior of the validity half-plane. The factor \((s)_4\) vanishes there; its accompanying integral is holomorphic there, so the remainder is zero. The factor \((s)_3\), multiplying the term with \(B_4\), also vanishes. Only the following remains:

\[
Z(-1)=\sum_{n=1}^N n-\frac{N^2}{2}-\frac N2-\frac{B_2}{2}.
\]
Since additive accumulation gives \(\sum_{n=1}^N n=N(N+1)/2\), and the Bernoulli recurrence gives \(B_2=1/6\), the result is \(-1/12\), independently of \(N\). \hfill\(\square\)

The evaluated object is the continuation of the functional, not the ordinary limit of the positive partial sums, which grow without bound. This distinction preserves the operation giving rise to the value: the Bernoulli term remains after cancellation of the polynomial boundary terms. Regularization therefore has a structure and genealogy here, in addition to a value.

\subsection{Laurent coefficients and preservation of the three classes}
Write

\[
Z(1+z)-\frac1z=\sum_{n\geq0}a_nz^n,
\qquad \gamma_n=(-1)^nn!a_n.
\]
The regular function on the left is holomorphic around zero. For a circle contained in its domain of holomorphy, Cauchy's formula determines

\[
a_n=\frac1{2\pi i}\int_{|z|=\rho}
\frac{Z(1+z)-1/z}{z^{n+1}}\,dz.
\]
The continuation formula with remainder controls this extraction: a uniform error bound \(\varepsilon\) on the circle produces a bound \(\varepsilon\rho^{-n}\) for the coefficient. This justifies a procedure at arbitrary order; it does not turn agreement of two finite quadratures into an error bound.

The residual readout maintains the three classes of the integer operator:

\[
Q_re_n=\mathbf1_{n\equiv r\ (3)}e_n,\qquad
Z_r(s)=\operatorname{Tr}(Q^{-s}Q_r),\quad r=0,1,2.
\]
These projections have a common domain in the convergence half-plane. Their sum is the identity; the difference of classes 1 and 2 preserves the orientation of the character modulo three. One has

\[
Z=Z_0+Z_1+Z_2,\qquad Z_0(s)=3^{-s}Z(s),\qquad
L_{-3}(s)=Z_1(s)-Z_2(s).
\]
The last series has periodic coefficients \((1,-1,0)\) and zero mean. Grouping by periods allows continuation near \(s=1\). Each \(Z_r\) has residue \(1/3\) there.

Let

\[
Z_r(1+z)-\frac1{3z}=\sum_{n\geq0}c_{r,n}z^n,
\qquad \lambda=\log3.
\]
Coefficient comparison gives

\[
a_n=c_{0,n}+c_{1,n}+c_{2,n},\qquad
\frac{L_{-3}^{(n)}(1)}{n!}=c_{1,n}-c_{2,n},
\]
\[
c_{0,n}=\frac13\left[
\frac{(-\lambda)^{n+1}}{(n+1)!}
+\sum_{k=0}^na_k\frac{(-\lambda)^{n-k}}{(n-k)!}\right].
\]
The third identity comes from multiplying the complete series, including its pole, by \(3^{-1}e^{-\lambda z}\). Sum and difference uniquely recover the other two coefficients. The three readouts preserve aggregation, orientation and change of scale, respectively.

At order zero, evaluation can be completed without receiving a decimal value of \(L_{-3}(1)\). Indeed,

\[
L_{-3}(1)=\sum_{m\geq0}\left(\frac1{3m+1}-\frac1{3m+2}\right)
=\int_0^1\frac{1-x}{1-x^3}\,dx
=\int_0^1\frac{dx}{1+x+x^2}.
\]
Complete the square in the denominator and evaluate the antiderivative:

\[
L_{-3}(1)=\frac2{\sqrt3}
\left[\arctan\frac{2x+1}{\sqrt3}\right]_0^1
=\frac{\pi}{3\sqrt3}.
\]
The coordinate \(\pi\) is HMT's prior regional representation; this identity is a subsequent analytic evaluation of a class readout, not a procedure for retrospectively selecting that coordinate. Thus, with \(C_r=c_{r,0}\),

\[
C_0=\frac{\gamma_0-\log3}{3},\qquad
C_{1,2}=\frac{\gamma_0}{3}+\frac{\log3}{6}
\mathbin{\pm}\frac{\pi}{6\sqrt3}.
\]
The following identities are recovered:

\[
C_0+C_1+C_2=\gamma_0,\qquad
C_1-C_2=L_{-3}(1),\qquad C_1+C_2-2C_0=\log3.
\]
\clearpage
\subsection{Exact representation and scope of the certificate}
The program \texttt{partes\_\allowbreak{}finitas\_\allowbreak{}exactas.\allowbreak{}py} preserves the following order: word sequence through the additive sheet; evaluation of integer modes; rational construction of Bernoulli numbers; application of the functional's formula; incorporation of the remainder; directed decimal representation. Native multiplication remains available in the same package and its compatibility is proved in the product chapter. Analytic computation on the expressed modes does not modify their rules.

With \(N=81\), \(p=8\) and 120 terms for each logarithm series, the resulting intervals are

\[
\begin{gathered}
0.577215664901532860606512090082273204041535\\
<\gamma_0<\\
0.577215664901532860606512090082531387279784,
\end{gathered}
\]
\[
\begin{gathered}
1.202056903159594285399738161511447311355829\\
<Z(3)<\\
1.202056903159594285399738161511452663119342.
\end{gathered}
\]
The decimal endpoints are outward roundings of exact fractions. The program also preserves each interval's rational width and recovers \(-1/12\) exactly. No comparison digit appears in the producer. The separate checker opens reference values after obtaining the outputs; this comparison checks execution and does not replace the proved bounds.

Within its domain, this chapter completes the finite parts and evaluations it has just defined. It uses no assumption about the location of zeros or positivity of the Weil form. Prolongation towards the spectral problem preserves the obligations of its own operators: a regularized evaluation does not by itself provide the boundary identity of another functional.

Provenance: development recovered from the triadic functionals and coordinated readout network. Bringing together the Euler–Maclaurin proof, its rational bounds and its program constitutes a formalization and additional certificate of this edition; no new historical priority is claimed for those analytic identities.
